\scanpage{16}
\diagram{\begin{tikzcd}[column sep=small]
\operatorname{Autext}_{loc}(\widehat\Pi_{g,\nu})\arrow[d,phantom,"\cup"']&\text{de }\operatorname{Isoext}_{loc}(\widehat\Pi_{g,\nu},\Pi_{X,x}),\ \text{au ss-groupe}\\
\mathcal N'_{g,\nu}&\text{image de }\mathcal N_{g,\nu}
\end{tikzcd}}
Mais si $\mathcal N_{g,\nu}\xrightarrow{\mathrm{épi}}\mathcal N'_{g,\nu}$ n'était un iso, il serait quand même plus fin de dire que via (13), on a un torseur can
\diagram{\begin{tikzcd}
(15)\quad P_{X,x}\arrow[d,equal]&\text{torseur à droite sous }\mathcal M_{g,\nu}\\
\operatorname{Isom}_{\substack{\text{groupoïde}\\\text{fond. de}\\\unsure{\mathcal X}_{g,\nu\mathbb Q}}}(x_0,x)&
\end{tikzcd}}
et un hom
\diagram{\begin{tikzcd}
(16)\quad P_{X,x}\arrow[r]&\operatorname{Isom}_{loc}(\widehat\Pi_{g,\nu},\Pi_{X,x})
\end{tikzcd}}
compatible avec
\diagram{\begin{tikzcd}
\deleted{(17)}\quad\mathcal M_{g,\nu}\arrow[r,"can"]&\operatorname{Aut}\deleted{\mathrm{ext}}_{loc}(\widehat\Pi_{g,\nu})
\end{tikzcd}}
sur les groupes d'opérateurs, d'où par passage aux quotients par $\widehat\Pi_{g,\nu}$
\diagram{\begin{tikzcd}
(17)\quad Q_{X\deleted{\unsure{,x}}}\arrow[r]\arrow[d,equal]&\operatorname{Isomext}_{loc}(\widehat\Pi_{g,\nu},\Pi_{X,x})\\
(P_{X,x}/\widehat\Pi_{X,x}\simeq\operatorname{Isom}_{\substack{\text{groupoïde}\\\text{fond. de }M_{g,\nu\mathbb Q}}}(\xi_0,\xi))&
\end{tikzcd}}
compatible avec
\diagram{\begin{tikzcd}
\mathcal N_{g,\nu}\arrow[r]&\operatorname{Autext}_{loc}(\widehat\Pi_{g,\nu},\Pi_{X,x})
\end{tikzcd}}
\deleted{Et (17)} Les données (16) et (17) sont équivalentes. (en tant que structures suppl. sur $\Pi_{X,x}$) -- c'est une ``arithmétisation'' en un sens plus fin que via \deleted{\unsure{un}} $P'_{X,x}$ ou $Q'_X$ (image de $P_{X,x}$, $Q_X$)

\mspage{15}
